\documentclass[10pt]{article}
\usepackage[UTF8]{ctex}
\usepackage{amsmath, amssymb}
\usepackage{geometry}
\geometry{a4paper, margin=1.5cm}
\usepackage{enumitem}

\title{计数原理应用教程：从题目到模型}
\author{}
\date{}

\begin{document}
\maketitle

\section{导言：如何阅读一道计数题？}

拿到一道计数题，请依次问自己三个问题：

\begin{enumerate}
\item \textbf{这件事是“分类”完成还是“分步”完成？}
\begin{itemize}
\item 分类：结果属于互斥的几类之一 → 加法原理
\item 分步：需要依次做几个动作 → 乘法原理
\item 多数综合题：先分类，再对每类分步；或先分步，其中某一步再分类。
\end{itemize}
\item \textbf{选出的元素是否考虑顺序？}
\begin{itemize}
\item 顺序重要 → 排列 \(A_n^m\)
\item 顺序不重要 → 组合 \(C_n^m\)
\item 注意：有时题目语言不直接，需要自己判断（例如“安排职务”有顺序，“选出代表”无顺序）。
\end{itemize}
\item \textbf{是否有特殊限制条件？}
\begin{itemize}
\item 特殊元素（如“0不能排在首位”）
\item 特殊位置（如“甲不能站在排头”）
\item 相邻/不相邻
\item 至多/至少
\item 分组分配（均匀/不均匀）
\end{itemize}
\end{enumerate}

下面按高考常见题型分类讲解。

\section{题型一：数字组数问题}

\textbf{典型题目}：用数字0,1,2,3,4组成无重复数字的三位数，共有多少个？

\subsection*{思考历程}
\begin{itemize}
\item \textbf{目标}：组成一个三位数（百位、十位、个位）。
\item \textbf{限制}：数字不重复，且百位不能为0。
\item \textbf{模型判断}：顺序显然重要（百十个位不同，数字不同），且选3个数字排列 → 排列。但因为有“0不能做百位”的特殊元素，不能直接 \(A_5^3\)。
\end{itemize}

\textbf{解法1（位置优先）}：
\begin{itemize}
\item 先安排百位：从\{1,2,3,4\}中选1个 → 4种。
\item 再安排十位：从剩下的4个数字（含0）中选1个 → 4种。
\item 最后个位：从剩下的3个中选1个 → 3种。
\end{itemize}
总数：\(4 \times 4 \times 3 = 48\)。

\textbf{解法2（元素优先，间接法）}：
\begin{itemize}
\item 不考虑0的限制，\(A_5^3 = 5\times4\times3=60\)。
\item 减去“0在百位”的情况：固定百位为0，再从剩余4个数字选2个排十位和个位 → \(A_4^2 = 4\times3=12\)。
\end{itemize}
总数：\(60-12=48\)。

\begin{quote}
\textbf{二级结论}：含0的数字组数问题，常用“先排特殊位置”或“总数减非法”。
\end{quote}

\subsection*{易错点}
\begin{itemize}
\item 忘记0不能做首位，直接 \(A_5^3\)。
\item 分步时，百位选完后，剩余数字个数要考虑是否包含0，但方法数不受影响（因为0没有被排除）。
\end{itemize}

\section{题型二：排队问题（人/元素排成一列）}

\textbf{典型题目}：5个人站成一排，求下列情况下的排法数：
\begin{enumerate}
\item 甲不站在排头；
\item 甲、乙相邻；
\item 甲、乙不相邻；
\item 甲在乙的左边（不一定相邻）。
\end{enumerate}

\subsection*{思考历程}

\subsubsection*{(1) 甲不站排头}
\textbf{方法一（位置优先）}：排头有4种选择（除甲外），剩下4个位置全排列 \(4!\) → \(4 \times 4! = 96\)。

\textbf{方法二（元素优先，间接法）}：总排法 \(5! = 120\)，减去甲在排头的情况：固定甲在排头，其余4人全排列 \(4! = 24\) → \(120-24=96\)。

\subsubsection*{(2) 甲、乙相邻}
\textbf{捆绑法}：将甲、乙看作一个整体（内部可交换顺序：2种）。现在共有：这个整体 + 其余3人 = 4个元素，全排列 \(4!\)。总数：\(2 \times 4! = 48\)。

\begin{quote}
\textbf{二级结论}：相邻问题用捆绑法，注意内部排列。
\end{quote}

\subsubsection*{(3) 甲、乙不相邻}
\textbf{方法一（间接法）}：总数 - 相邻数 = \(120 - 48 = 72\)。

\textbf{方法二（插空法）}：先排其余3人，有 \(3! = 6\) 种排法，产生4个空（包括两端）。从4个空中选2个插入甲、乙（顺序重要：\(A_4^2 = 4\times3=12\)）。总数：\(6 \times 12 = 72\)。

\begin{quote}
\textbf{二级结论}：不相邻问题用插空法，先排无限制元素，再在空位中插入特殊元素。
\end{quote}

\subsubsection*{(4) 甲在乙的左边（不一定相邻）}
\textbf{对称法}：在全部排列中，甲在乙左边和甲在乙右边的排列数相等，且无其他可能。所以甲在乙左边的排法数为 \(\frac{5!}{2} = 60\)。

\begin{quote}
\textbf{二级结论}：对于两个不同元素，要求顺序固定（如左/右、前/后），排法数为总排列数除以2。
\end{quote}

\subsection*{易错点}
\begin{itemize}
\item 捆绑法忘记内部排列（如把甲、乙捆成一个，但内部有2种顺序）。
\item 插空法：空位数计算错误（n个人有n+1个空）。
\item 对称法：只有当两个元素地位完全对称且无其他限制时才适用。
\end{itemize}

\section{题型三：分配问题（不同元素分给不同对象）}

\textbf{典型题目}：将4本不同的书分给3个人，每人至少1本，有多少种分法？

\subsection*{思考历程}
\begin{itemize}
\item \textbf{问题本质}：把4个不同元素分成3组（有标签，因为人是不同的），且每组非空，再将组分配给不同人。
\item 因为书不同，人不同，所以是\textbf{先分组（组合），再分配（排列）}。
\end{itemize}

\textbf{步骤}：
\begin{enumerate}
\item 将4本不同的书分成3堆（无序的组），其中一堆有2本，另外两堆各1本。
\begin{itemize}
\item 选2本书作为一堆：\(C_4^2 = 6\)。
\item 剩下2本书各自成堆。更严谨的做法：分组方式为“2,1,1”。对于不同元素，分组数为 \(\frac{C_4^2 \cdot C_2^1 \cdot C_1^1}{2!}\)，因为两个1本组对称。计算：\(\frac{6 \times 2 \times 1}{2} = 6\)。
\end{itemize}
\item 将分好的3组（无序）分配给3个不同的人：有 \(3! = 6\) 种分配方式。
\end{enumerate}
总数：\(6 \times 6 = 36\)。

\begin{quote}
\textbf{二级结论}：不同元素分配给不同对象（每人至少一个），常用“先分组再分配”。分组时要留意是否有相同大小的组，需要除以对称性阶乘。
\end{quote}

\subsection*{易错点}
\begin{itemize}
\item 直接 \(4^3\)（允许有人没拿到书），但题目要求每人至少1本。
\item 分组时忘记除以对称阶乘，导致重复计数。
\end{itemize}

\section{题型四：涂色问题}

\textbf{典型题目}：用5种不同颜色给如图的4个区域涂色，要求相邻区域颜色不同，有多少种涂法？（以一条直线上的4个区域为例，区域1-2-3-4依次相邻）

\subsection*{思考历程}
\begin{itemize}
\item \textbf{分步涂色}：
\begin{itemize}
\item 区域1：5种。
\item 区域2：不能与1同色，4种。
\item 区域3：不能与2同色，但可以与1同色（因为不相邻），所以有4种（除去区域2的颜色）。
\item 区域4：不能与3同色，有4种（除去区域3的颜色）。
\end{itemize}
总数：\(5 \times 4 \times 4 \times 4 = 320\)。
\end{itemize}

\textbf{变式}：若区域首尾也相邻（环形），则需更复杂处理，通常分情况讨论。

\begin{quote}
\textbf{二级结论}：直线型相邻涂色，除了第一步，后面每步都是（颜色总数-1），但要警惕区域是否与前面多个区域相邻（如格子图需用“分类”）。
\end{quote}

\subsection*{易错点}
\begin{itemize}
\item 误以为区域3只能从剩余3种颜色中选（忽略了它可以与区域1同色）。
\item 环形涂色常用“先选一种颜色给某区域，再分类讨论”。
\end{itemize}

\section{题型五：至多、至少问题（用组合数的性质）}

\textbf{典型题目}：从6名男生和4名女生中选出4人，要求至少有2名女生，有多少种选法？

\subsection*{思考历程}
\begin{itemize}
\item \textbf{直接分类}：
\begin{itemize}
\item 女生2人+男生2人：\(C_4^2 \cdot C_6^2\)
\item 女生3人+男生1人：\(C_4^3 \cdot C_6^1\)
\item 女生4人：\(C_4^4\)
\end{itemize}
总数 = 这些之和。
\item \textbf{间接法}：
\begin{itemize}
\item 总选法 \(C_{10}^4\) 减去不符合条件（女生0人或1人）
\item 女生0人：\(C_6^4\)
\item 女生1人：\(C_4^1 \cdot C_6^3\)
\end{itemize}
总数 = \(C_{10}^4 - C_6^4 - C_4^1 \cdot C_6^3\)。
\end{itemize}

\begin{quote}
\textbf{二级结论}：“至少”问题常用两种方法：直接分类（注意不重不漏）或总数减反面的“至多”情况。
\end{quote}

\subsection*{易错点}
\begin{itemize}
\item 分类时漏掉某种情况（如至少2人包括2,3,4）。
\item 间接法时反面情况计算错误。
\end{itemize}

\section{题型六：分组分配中的均匀分组问题}

\textbf{典型题目}：将6本不同的书分成3堆，每堆2本，有多少种分法？

\subsection*{思考历程}
\begin{itemize}
\item 这是一个\textbf{均匀分组}（无标签）问题。
\item 步骤：
\begin{enumerate}
\item 先依次取出：第一堆 \(C_6^2\)，第二堆 \(C_4^2\)，第三堆 \(C_2^2\)，得到 \(C_6^2 C_4^2 C_2^2 = 15 \times 6 \times 1 = 90\)。
\item 但这样给三堆赋予了顺序（第一、第二、第三），而实际上三堆是无区别的。每3!种排列对应同一种分组。所以实际分法为 \(\frac{90}{3!} = 15\)。
\end{enumerate}
\end{itemize}

\begin{quote}
\textbf{二级结论}：均匀分组（各组元素个数相同）且组无标签时，需除以组数的阶乘。若组有标签（如分给三个不同的人），则不需要除以阶乘，直接 \(C_6^2 C_4^2 C_2^2\) 即可。
\end{quote}

\subsection*{易错点}
\begin{itemize}
\item 混淆“分组”与“分配”：有标签的分配不除阶乘，无标签的分组要除。
\item 部分均匀分组（如2,2,1,1）也要除对称组的阶乘。
\end{itemize}

\section{题型七：隔板法（相同元素分配问题）}

\textbf{典型题目}：将10个相同的球放入3个不同的盒子，每个盒子至少1个，有多少种放法？

\subsection*{思考历程}
\begin{itemize}
\item 问题本质：把10个相同元素分成3组（非空），组有标签（盒子不同）。
\item 用\textbf{隔板法}：10个球排成一排，中间有9个空隙。要分成3组，需插入2块隔板，且不能放在两端（因为非空）。从9个空隙中选2个放隔板：\(C_9^2 = 36\)。
\end{itemize}

\textbf{推广}：若允许空盒，则等价于方程 \(x_1+x_2+x_3=10\)，\(x_i \ge 0\)，非负整数解个数为 \(C_{10+3-1}^{3-1}=C_{12}^2\)。

\begin{quote}
\textbf{二级结论}：相同元素分给不同对象（允许空）用隔板法：\(C_{n+k-1}^{k-1}\)，其中n为元素个数，k为对象数。
\end{quote}

\subsection*{易错点}
\begin{itemize}
\item 隔板法仅适用于相同元素。
\item 非空与允许空时，公式不同。
\end{itemize}

\section{总结：解题的一般流程}

\begin{enumerate}
\item \textbf{读题，划出关键信息}：元素是否相同？对象是否不同？是否允许重复？是否有顺序？有无特殊限制？
\item \textbf{判断模型}：
\begin{itemize}
\item 若元素不同且考虑顺序 → 排列（或分步乘法）
\item 若元素不同且不考虑顺序 → 组合
\item 若元素相同 → 隔板法或枚举
\item 若既有分类又有分步 → 先分类，每类中分步
\end{itemize}
\item \textbf{选择策略}：直接法、间接法、捆绑法、插空法、先分组再分配、隔板法等。
\item \textbf{检查}：是否遗漏或重复？特殊条件是否满足？数字是否合理（如概率题中分母是否一致）？
\end{enumerate}

\section*{最后的话}

计数原理的题目千变万化，但核心思想只有两个：\textbf{加法原理}与\textbf{乘法原理}。所有的排列组合公式、二级结论，都是这两个原理在不同场景下的应用。

当你遇到陌生题目时，不要急于套用公式，而是回到基本原理，自己设计“分步”或“分类”的方案。真正重要的不是记住多少种题型，而是培养\textbf{将实际问题转化为计数模型}的能力。

\begin{quote}
\textbf{后续练习建议}：每做一道题，在草稿纸上写下你的“思考历程”三问（分类还是分步？顺序是否重要？特殊条件如何处理？），并尝试用至少两种方法验证答案。
\end{quote}

\end{document}